% English translation of chapter3.tex lines 9--55.
% Source: Methods of Algebra, Volume 2, commit
% 9a5803ff77dd3257484cb177f851a73770a59dd3 (CC BY 4.0).
% stable-unit-id: o014.aljabr2.chapter3.overview
% Scope: Chapter 3 overview before Section 3.1.

% segment-id: o014.li2.chapter3.overview.h001
\chapter{Complexes}\label{sec:cplx}
\hypertarget{o014.aljabr2.chapter3.overview}{ }

% segment-id: o014.li2.chapter3.overview.p001
The notion of a complex was introduced briefly in \S~\ref{sec:cohomology}.
For a complex $X=(X^n,d^n)_n$ in an abelian category $\mathcal{A}$, the
cohomology
$\Hm^n(X):=\Ker(d^n)/\Image(d^{n-1})$%
\footnote{Editorial correction (O014-C032): the source prints
$\Image(d^{n+1})$. Under the cochain convention
$X^{n-1}\xrightarrow{d^{n-1}}X^n\xrightarrow{d^n}X^{n+1}$, the subobject
contained in $\Ker(d^n)$ is $\Image(d^{n-1})$; the corrected index agrees
with the book's earlier definition.}
extracts important invariants from $X$. Conversely,%
\footnote{Editorial correction (O014-C030): the source has a full stop
followed by a stray character before the clause about the vanishing of all
cohomology, leaving the transition ungrammatical. The context gives the
converse direction, restored here without changing the mathematical claim.}
a complex whose cohomology vanishes in every degree is an exact sequence,
also called an acyclic complex.

% segment-id: o014.li2.chapter3.overview.p002
Much of \S\S~\sourcecrossref{sec:Hom-cplx}{3.2}--%
\sourcecrossref{sec:double-cplx}{3.5} does not involve cohomology and applies
to an arbitrary additive category $\mathcal{A}$. The category of complexes
over $\mathcal{A}$ is denoted by $\cate{C}(\mathcal{A})$. Unless stated
otherwise, all complexes in this chapter are cochain complexes in the sense
of Remark~\ref{rem:cochain-vs-chain}; whenever $\Bbbk$ occurs, it denotes an
arbitrary commutative ring.

% segment-id: o014.li2.chapter3.overview.p003
For complexes $X$ and $Y$ over an additive category $\mathcal{A}$, homotopy
is an equivalence relation on
$\Hom_{\cate{C}(\mathcal{A})}(X,Y)$. It can also be understood through the
$\Hom$ complex $\Hom^\bullet(X,Y)$, studied in
\S~\sourcecrossref{sec:Hom-cplx}{3.2}. If $\mathcal{A}$ is abelian,
homotopic morphisms induce the same maps on cohomology. The mapping cone
$\Cone(f)$ introduced in \S~\sourcecrossref{sec:mapping-cone}{3.3} will
subsequently play a major role in the study of derived categories. Mapping
cones are closely related to homotopy, and both have direct topological
interpretations.

% segment-id: o014.li2.chapter3.overview.p004
A double complex
$X=(X^{p,q},\dhori^{p,q},\dvert^{p,q})_{p,q}$ may be viewed as a complex
with two independent directions, or equivalently as a ``complex of
complexes.'' A double complex $X$ can be flattened by direct sums
(respectively, direct products) into a complex $\tot_{\oplus}X$
(respectively, $\tot_{\Pi}X$), called a total complex of $X$. The precise
definitions involve several signs; these details are the subject of
\S~\sourcecrossref{sec:double-cplx}{3.5}.

% segment-id: o014.li2.chapter3.overview.p005
One of the basic techniques of homological algebra is to derive a long exact
sequence in cohomology from a short exact sequence of complexes
$0\to X'\xrightarrow{f}X\xrightarrow{g}X''\to0$:
% segment-id: o014.li2.chapter3.overview.q001
\[
  \cdots \to \Hm^n(X') \xrightarrow{\Hm^n(f)} \Hm^n(X)
  \xrightarrow{\Hm^n(g)} \Hm^n(X'')
  \xrightarrow{\text{connecting morphism}} \Hm^{n+1}(X') \to \cdots .
\]
Long exact sequences can also be understood from the viewpoint of mapping
cones, although this requires a nontrivial argument. These topics are treated
in \S\S~\sourcecrossref{sec:Abel-cplx}{3.6}--%
\sourcecrossref{sec:cone-vs-long-exact-sequence}{3.7}.

% segment-id: o014.li2.chapter3.overview.p006
To help the reader become familiar with the basic operations on complexes,
\S~\sourcecrossref{sec:HH}{3.8} offers an integrated exercise based on the
Hochschild homology $\HHm_n(M)$ and Hochschild cohomology $\HHm^n(M)$ of an
$(R,R)$-bimodule $M$. Hochschild homology and cohomology have many uses
throughout mathematics. This chapter gives only a brief pedagogical
introduction; both will reappear as examples in later chapters.

% segment-id: o014.li2.chapter3.overview.p007
Using the truncation functors introduced in
\S~\sourcecrossref{sec:truncation-functors}{3.9}, we shall relate the
vertical or horizontal cohomology of a double complex to the cohomology of
its total complex in \S~\sourcecrossref{sec:double-cplx-coh}{3.10}. Certain
boundedness conditions on the double complex are needed. Some textbooks
treat these properties with spectral sequences; following \cite{KS06}, this
book takes a somewhat indirect route and proves them directly.

% segment-id: o014.li2.chapter3.overview.p008
For an object $X$ of an abelian category $\mathcal{A}$, a resolution of $X$
has left and right versions, concretely given by exact sequences%
\footnote{Editorial correction (O014-C031): the source prints
$0\to X\to I^0\to I^2\to\cdots$, omitting $I^1$. The diagram immediately
following uses $I^0,I^1,I^2$ in order; the missing term $I^1$ is restored here.}
% segment-id: o014.li2.chapter3.overview.q002
\[
  0\to X\to I^0\to I^1\to I^2\to\cdots
  \quad\text{or}\quad
  \cdots\to P_1\to P_0\to X\to0 .
\]
Equivalently, one replaces $X$ by a complex with better properties through a
quasi-isomorphism $X\to I:=(I^n)_n$ or $P:=(P_n)_n\to X$, where $P$ is a
chain complex. A quasi-isomorphism is a morphism inducing isomorphisms on
cohomology or homology. If every $I^n$ (respectively, $P_n$) is required to
be injective (respectively, projective), the construction is called an
injective (respectively, projective) resolution of $X$. Resolutions are
studied in depth in \S~\sourcecrossref{sec:resolutions}{3.11}. They are also
functorial and unique up to homotopy in the following sense. For the
injective version, suppose the solid part of the following diagram lies in
$\mathcal{A}$.
% segment-id: o014.li2.chapter3.overview.g001
\[
\begin{tikzcd}
	0 \arrow[r] & X' \arrow[r] & (I')^0 \arrow[r] & (I')^1 \arrow[r]
	  & (I')^2 \arrow[r] & \cdots \\
	0 \arrow[r] & X \arrow[u,"f"] \arrow[r]
	  & I^0 \arrow[dashed,u,"{\beta^0}"] \arrow[r]
	  & I^1 \arrow[r] \arrow[dashed,u,"{\beta^1}"]
	  & I^2 \arrow[r] \arrow[dashed,u,"{\beta^2}"] & \cdots
\end{tikzcd}
\]
Assume that both rows are exact and that every $(I')^n$ is injective. There
are dashed arrows $\beta^0,\beta^1,\ldots$ making the entire diagram
commutative, and the morphism of complexes $\beta:I\to I'$ is unique up to
homotopy. This is a joint application of
Theorem~\sourcecrossref{prop:resolution-homotopy}{in \S~3.11} and
Lemma~\sourcecrossref{prop:ext-alpha-beta}{in \S~3.11}, both special cases of
broader statements.

% segment-id: o014.li2.chapter3.overview.p009
In fact, \S~\sourcecrossref{sec:resolutions}{3.11} treats not only
resolutions of objects $X$ of $\mathcal{A}$, but also resolutions of
complexes. Its results include the useful Cartan--Eilenberg resolution
(Theorem~\sourcecrossref{prop:CE-resolution}{in \S~3.11}). Some arguments for
complexes inevitably become lengthy, even though their underlying idea
remains clear.

% segment-id: o014.li2.chapter3.overview.p010
Replacing objects of $\mathcal{A}$ by injective (respectively, projective)
resolutions, applying a left exact (respectively, right exact) functor
$F:\mathcal{A}\to\mathcal{B}$ termwise, and then taking $\Hm^n$
(respectively, $\Hm_n=\Hm^{-n}$) leads to the classical definition of the
right derived functors $\mathrm{R}^nF$ (respectively, the left derived
functors $\mathrm{L}_nF$). Short exact sequences in $\mathcal{A}$ then induce
long exact sequences of derived functors; see
\S~\sourcecrossref{sec:derived-primer}{3.12}. Resolutions of complexes also
allow derived functors to be evaluated on bounded-below or bounded-above
complexes; in the classical literature these constructions are called
hyperderived functors.

% segment-id: o014.li2.chapter3.overview.p011
Right (respectively, left) derived functors and their long exact sequences
naturally give rise to cohomological (respectively, homological)
$\delta$-functors, among which derived functors are ``universal.''
Proposition~\sourcecrossref{prop:erasable-univ-delta}{in \S~3.12}
characterizes universal $\delta$-functors by the erasability (respectively,
co-erasability) condition of Definition~\sourcecrossref{def:erasable}{in
\S~3.12}. This result is due to Grothendieck. The associated collection of
techniques forms the core of classical homological algebra.

% segment-id: o014.li2.chapter3.overview.p012
Another important construction consists of the right (respectively, left)
derived functors of a bifunctor
$F:\mathcal{A}_1\times\mathcal{A}_2\to\mathcal{B}$ such that $F$ is left
exact (respectively, right exact) in both variables. Standard examples are the
$\Ext$ and $\Tor$ functors of \S~\sourcecrossref{sec:Ext-Tor}{3.14}; the
former necessarily involves complexes over an opposite category.
Theorem~\sourcecrossref{prop:balanced-primer}{in \S~3.14} studies a class of
bifunctors called balanced and shows that deriving in either variable gives
the same result. Its proof rests on basic properties of the cohomology of
double complexes.

% segment-id: o014.li2.chapter3.overview.p013
Suppose that $\mathcal{A}$ has exact countable products. The construction
$\lim^1$ discussed in \S~\sourcecrossref{sec:lim1}{3.13} is another example
of a right derived functor: it is the first right derived functor of
$\varprojlim$, while all higher right derived functors vanish
(Theorem~\sourcecrossref{prop:lim1}{in \S~3.13}). To obtain practical
criteria ensuring $\lim^1=0$, we shall introduce the Mittag--Leffler
condition in Definition~\sourcecrossref{def:ML}{in \S~3.13}.

% segment-id: o014.li2.chapter3.overview.p014
Finally, in \S~\sourcecrossref{sec:K-injectives}{3.15} we study
K-injective and K-projective complexes and define K-injective and
K-projective resolutions. These are natural generalizations of injective and
projective resolutions to unbounded complexes. Because K-injectivity and
K-projectivity are properties of an entire complex, their definitions are
more concise and natural than the termwise definitions of injective and
projective resolutions. They extend the definition of derived functors to
unbounded complexes. The issue is existence: Theorems
\sourcecrossref{prop:K-injective-resolution}{in \S~3.15} and
\sourcecrossref{prop:K-projective-resolution}{in \S~3.15} give partial
results. Their proofs are comparatively intricate and require some
properties of $\lim^1$.

% segment-id: o014.li2.chapter3.overview.n001
\begin{readingnote}
	Experience with complexes of $\Bbbk$-modules, where $\Bbbk$ is a
	commutative ring, helps one understand this chapter more quickly but is not
	strictly necessary. After Chapter~\ref{sec:Abel-cat}, working with abelian
	categories more general than module categories presents no fundamental
	difficulty.

	Readers who want to reach the classical definition of derived functors
	quickly and are not interested in its topological motivation may first skip
	the material on mapping cones and mapping cylinders in
	\S~\sourcecrossref{sec:mapping-cone}{3.3},
	\S~\sourcecrossref{sec:opposite-cplx}{3.4}, and
	\S~\sourcecrossref{sec:cone-vs-long-exact-sequence}{3.7}. In addition,
	\S~\sourcecrossref{sec:opposite-cplx}{3.4} relates
	$\cate{C}(\mathcal{A}^{\opp})$ to $\cate{C}(\mathcal{A})^{\opp}$. Some of
	the signs there require care, but they do not greatly affect the main ideas.
	On a first reading, one may study only the statements and skip the proofs.

	Since the later material does not depend on Hochschild homology or
	cohomology, readers may decide whether to skip it according to their time
	and interests. It is not difficult and will recur as examples and remarks
	in subsequent chapters.

	Similarly, the material of \S~\sourcecrossref{sec:lim1}{3.13} appears only
	as examples later, except in \S~\sourcecrossref{sec:K-injectives}{3.15}.
	The material of \S~\sourcecrossref{sec:K-injectives}{3.15} is cited only
	in \S~\sourcecrossref{sec:unbounded-derived}{4.11} and
	\S~\sourcecrossref{sec:otimesL}{4.12}. Readers who are not particularly
	interested in unbounded derived functors may first study the definitions
	and Example~\sourcecrossref{eg:bdd-K-resolution}{in \S~3.15}, then omit
	the remainder.
\end{readingnote}
